\chapter{Deciding It}\label{ch:deciding}

Theorem~\ref{thm:faithful} says what to check. This chapter says how to check it fast
enough to use interactively on real protocols, and proves the costs. There are two
questions. The first is whether a candidate system is faithful to a protocol; it never
needs the formulas at all. The second is whether a model satisfies a CTL formula,
including the large formulas faithful rules and project rules can be; it needs an
evaluator whose cost does not grow with the formula's size as a tree.

\section{Faithfulness without formulas}

Checking the rules one at a time costs one model check per typestate, and each model
check is linear in the model: quadratic in all. The rules, taken together, only say
that the candidate is bisimilar to the protocol (Theorem~\ref{thm:faithful}), and
bisimilarity can be decided directly.

\paragraph{Refinement.} LeanFM computes the rounds of Definition~\ref{def:rounds} over
the disjoint union of the protocol's model and the candidate's, and the candidate is
faithful exactly when the two start nodes end in one class. A round numbers, through a
hash table, each pair (class, sorted set of successor classes): one pass over the
steps. The rounds stop when no class splits, after $K + 1$ rounds, $K$ the lookahead.

\begin{proposition}\label{prop:refine}
Refinement costs $O((K+1)(n + m\log d))$ hash-table operations for $n$ nodes, $m$ steps
and out-degree at most $d$, and $O(n \cdot m \log d)$ in the worst case.
\end{proposition}

\begin{proof}
Each round sorts each node's successor classes ($O(d\log d)$ per node, $O(m\log d)$ in
all) and looks up one key per node. There are $K+1$ rounds, and since each round that
does not stop splits at least one class, $K < n$.
\end{proof}

In every protocol LeanFM carries, $K \le 6$, so refinement is effectively linear. Paige
and Tarjan's algorithm achieves $O(m \log n)$ unconditionally
\citep{paige1987three}; LeanFM does not use it, and the worst case above is real for
protocols with long hidden counters.

\paragraph{The walk.} When the protocol's typestates are deterministic, LeanFM uses the
walk of Theorem~\ref{thm:walk} instead: one pass over the candidate's steps, $O(m)$.
This is the case prev makes possible.

\paragraph{When they differ.} For an unfaithful candidate, LeanFM searches breadth first
through pairs of nodes with equal labels that the refinement keeps apart, from the two
starts, for the first pair whose labels or next labels differ, and reports the path to
it. A candidate that never answers 404:

\begin{lstlisting}[language={},basicstyle=\ttfamily\footnotesize,breaklines=true]
Images: the candidate (3 nodes) is not faithful to the protocol (4 typestates), decided by walk: after GetImage, the protocol allows Image (may stop), NotFound (may stop) and the candidate Image (may stop)
\end{lstlisting}


A GoServer whose refused sign-in falls through to the room instead of retrying:

\begin{lstlisting}[language={},basicstyle=\ttfamily\footnotesize,breaklines=true]
GoServer: the candidate (39497 nodes) is not faithful to the protocol (1604 typestates), decided by refinement: after IssueSession -> SignInRefused, the protocol allows IssueSession and the candidate EnterRoom
\end{lstlisting}


And the classic pair that no trace comparison can separate, $a$ then $b$ or $a$ then
$c$ as two branches, against $a$ then a choice of $b$ or $c$: both produce exactly the
traces $ab$ and $ac$, yet they are not bisimilar, and CTL tells them apart
($\EX\,\AX\,\after{b}$ holds of one and not the other).

\begin{lstlisting}[language={},basicstyle=\ttfamily\footnotesize,breaklines=true]
faithful: false
method: refinement
path: [a]
protocol next: [b (may end)]
candidate next: [b (may end), c (may end)]
note:
\end{lstlisting}


\section{Evaluating a CTL formula}

A formula is evaluated by computing, for each subformula, the set of nodes where it
holds. LeanFM does this over a DAG, not a tree.

\paragraph{Hash-consing.} Each distinct subformula becomes one node of a DAG, found
through a hash table keyed by its operator, its children's node numbers and, for an
atom, the atom itself \citep{filliatre2006hashconsing}. The key is compared in full on
every hash match, so a collision costs a comparison and never merges two different
subformulas. A formula whose parts are already shared in memory, as the s-expression
reader builds them from labels, is walked once per shared part, so building the DAG is
linear in the DAG, not in the tree. Derived operators are rewritten into the core of
Definition~\ref{def:ctl} with the abbreviations given there.

\paragraph{Evaluation.} The DAG's nodes are numbered children first, and evaluated in
that order. Each truth set is a packed bitset over the model's nodes, and the model's
steps are held as offset and target arrays in both directions.
\begin{itemize}
\item $\neg$, $\wedge$, $\vee$ are word operations on bitsets.
\item $\mathsf{EX}\,\varphi$ and $\mathsf{AX}\,\varphi$ are one pass over the steps.
\item $\mathsf{E}[\varphi\,\mathsf{U}\,\psi]$ starts from the $\psi$-nodes and works
  backwards: a predecessor that satisfies $\varphi$ and is not yet in the set joins it
  and is pushed on a worklist.
\item $\mathsf{A}[\varphi\,\mathsf{U}\,\psi]$ keeps, per node, the number of its steps
  not yet known to lead into the set; when it reaches zero at a $\varphi$-node, the
  node joins.
\item $\mathsf{EG}\,\varphi$ starts from the $\varphi$-nodes and removes, by worklist,
  every node left with no successor inside.
\end{itemize}
Each truth set is freed as soon as every node that reads it has been evaluated, and the
bytes held at once are bounded; a check that would exceed the bound stops and says so,
and proves nothing.

\begin{theorem}\label{thm:eval}
Evaluating a formula with $D$ distinct subformulas over a model with $n$ nodes and $m$
steps takes $O(D \cdot (n + m))$ time.
\end{theorem}

\begin{proof}
Each DAG node is evaluated once. The Boolean operators and $\mathsf{EX}$, $\mathsf{AX}$
take $O(n + m)$. In the $\mathsf{EU}$ worklist a node is pushed at most once, when it
joins, and each push examines that node's predecessor steps once: $O(n + m)$. In
$\mathsf{AU}$ each step's counter is decremented at most once, when its target joins:
$O(n + m)$. In $\mathsf{EG}$ each node is removed at most once and each removal
examines its predecessor steps once: $O(n + m)$.
\end{proof}

These are the classical bounds of CTL model checking \citep{clarke1986automatic}, with
$D$ in place of the formula's length. The difference matters exactly where faithful
rules live: Retry's names are hundreds of nodes as trees and dozens as DAGs
(Chapter~\ref{ch:constructs}), and the replication model of
Chapter~\ref{ch:cases} has rules totalling about $6\times 10^9$ nodes as trees and
$14{,}757$ distinct subformulas. LeanFM's tests check the evaluator against the plain fixpoint
semantics at every node, for every formula of depth two over every operator, on a
draft with a loop, a choice, an ending and a deadlock and on every canonical task.

\section{Formulas that share}

A formula travels to LeanFM as JSON or as an s-expression. JSON is a tree and cannot
say ``this part again''; an s-expression can, with a label at the first occurrence and a
reference after it:
\begin{center}
\texttt{(and \#1=(EF (after c)) (AG (implies (after a) \#1\#)))}
\end{center}
The reader is data only: there is no read-time evaluation, and no \texttt{\#} form other
than \texttt{\#n=} and \texttt{\#n\#}. A label is defined only once its expression is
complete, so a reference inside its own expression is to an undefined label and is
refused; no cycle can be written. An answer whose formula would be too large to write
as a tree gives it back as an s-expression with its shared parts labelled.

\section{The protocols LeanFM carries}

\begin{table}[ht]
\centering
\begin{tabular}{@{}lrrrrl@{}}
\toprule
Protocol & Nodes & Typestates & Lookahead & Rules & Deterministic \\
\midrule
Ping (examples) & 4 & 4 & 0 & 5 & yes \\
PlaceOrder (examples) & 10 & 10 & 0 & 11 & yes \\
Ride (examples) & 28 & 16 & 0 & 17 & yes \\
BoundedRide (examples) & 45 & 23 & 6 & 24 & yes \\
UserLookup (examples) & 32 & 21 & 0 & 22 & yes \\
OAuth sign-in & 16 & 15 & 0 & 16 & yes \\
LeanFM's own MCP protocol & 101 & 33 & 1 & 34 & yes \\
IslandUploads Socket & 38 & 38 & 1 & 39 & yes \\
Maqam audio period & 1849 & 45 & 0 & 46 & yes \\
GoServer & 39497 & 1604 & 2 & 1605 & no \\
\bottomrule
\end{tabular}

\caption{Faithful CTL of the protocols in this repository: checker nodes, typestates,
lookahead, rules (one per typestate and one for the start) and whether the typestates
are deterministic.}
\label{tab:corpus}
\end{table}

Most protocols have lookahead $0$: the last message names the typestate and every rule
is the one-step shape. BoundedRide counts dispatch attempts and has lookahead $6$.
GoServer is the only one whose typestates are not deterministic; it is also the only
large one, and its size is the product of sessions that, checked one at a time, are
small (Proposition~\ref{prop:sessions}).

Measured through the MCP server on GoServer (39{,}497 checker nodes), on the author's
machine: all 1{,}605 faithful rules in 0.73\,s; GoServer against itself, by
refinement, in 1.42\,s; the mutant above in 1.42\,s; and each rule, sent back as an
s-expression and checked by \texttt{session\_ctl}, in about 0.84\,s, most of it building
the model. Before this evaluator, the same rules checked one at a time did not finish.
